\subsection{Sprint 0}

No cronograma do semestre, a apresentação da \ac{AGES} e a integração das equipes estavam previstas para 2 de março de 2026, e a apresentação do projeto pelos \textit{stakeholders}, para 9 de março. A apresentação das histórias de usuário e dos \textit{mockups} e o planejamento da \textit{Sprint} 1 estavam marcados para 25 de março, seguidos da retrospectiva e do relatório da \textit{Sprint} 0 em 30 de março.

A \textit{Sprint} 0 teve papel decisivo para o sucesso inicial do projeto, pois concentrou toda a etapa de estruturação conceitual, definição visual e preparação organizacional da equipe. Diferentemente de fases iniciais em que normalmente há menor intensidade técnica, nesta oportunidade o volume de trabalho foi elevado desde os primeiros encontros, especialmente pela necessidade de transformar uma ideia ampla em um produto viável e bem organizado.

Minha atuação concentrou-se principalmente na frente de produto, experiência do usuário e organização das entregas. Auxiliei diretamente colegas da \ac{AGES} I no desenvolvimento dos \textit{mockups}, esclarecendo dúvidas sobre boas práticas de interface, padronização visual e coerência entre telas. Busquei contribuir não apenas executando tarefas, mas também orientando membros menos experientes, compartilhando conhecimentos adquiridos em experiências anteriores e no ambiente profissional.

Participei ativamente da definição dos fluxos centrais da aplicação. Como o sistema possuía múltiplos perfis de usuários --- \ac{ONGs}, empresas investidoras e administradores --- foi necessário pensar cuidadosamente na jornada de cada tipo de usuário dentro da plataforma. Assim, colaborei na concepção lógica dos seguintes fluxos:

\begin{itemize}
\item cadastro institucional de \ac{ONGs};
\item cadastro empresarial;
\item portal administrativo;
\item listagem e gerenciamento de projetos;
\item sistema de vínculos entre organizações e empresas;
\item mural de editais;
\item denúncias de irregularidades;
\item edição de perfis institucionais;
\item visualização pública e privada de perfis.
\end{itemize}

Além disso, realizei \textit{benchmarking} de plataformas já consolidadas no mercado, analisando soluções semelhantes para identificar padrões modernos de navegação, \textit{layout}, \textit{onboarding} e \textit{dashboards}. Essa etapa foi importante para evitar decisões arbitrárias e elevar o nível de maturidade visual do produto.

Também participei da criação e refinamento de tarefas no Kanban, ajudando a transformar ideias amplas em demandas executáveis e priorizadas. Em diversos momentos, alinhei expectativas com colegas da \ac{AGES} IV sobre escopo, prioridades e viabilidade de determinadas funcionalidades.

Durante essa \textit{sprint}, percebi claramente a importância da comunicação em projetos colaborativos. Como havia muitas ideias simultâneas e diferentes níveis de experiência entre os integrantes, nem sempre o fluxo de informação acontecia da melhor forma. Isso exigiu postura mais ativa da minha parte, buscando centralizar dúvidas, destravar tarefas e orientar colegas.

Considero que a \textit{Sprint} 0 foi extremamente produtiva para mim, pois atuei em áreas diversas: \ac{UX}, produto, organização ágil e apoio técnico ao time. Também serviu para consolidar uma postura mais madura de liderança informal, auxiliando colegas e pensando no projeto de forma sistêmica, não apenas individual.

\vspace{0.5cm}
